\chapter{Experimental design}
\label{ch:design}

This chapter is the pre-registration of the final grid.
Everything in it was fixed before any run of that grid was launched: the
questions, the objects that answer them, the arms, the hypotheses with
their readouts and their nulls, the hyperparameters and what is checked
about them, the protocol, and the rules that decide what happens when a
check fails.
Chapter~\ref{ch:results} reports against this chapter.
The runs executed while the design was being fixed, fifteen epochs each,
are pilots; Section~\ref{sec:pilots} lists the design decisions they
settled and Section~\ref{sec:pilot-results} reports them as pilots.

\section{Questions and the shape of the evidence}
\label{sec:questions}

Three questions, in the order the evidence is built.
\begin{enumerate}
\item \textbf{Structure.} What does the exact solution of the game
  (Chapter~\ref{ch:env}) imply about what a prior-dominated learner can
  learn in it, and in particular whether a shaping objective has anything
  to buy?
  This is answered by dynamic programming before any training run, and it
  is the reason the thesis moves from the deterministic game to
  stochastic regeneration (Section~\ref{sec:a-to-b}).
\item \textbf{Shaping.} Does the trial-level objective of
  \eqref{eq:shaper-objective} change what a naive learner acquires, and
  what the shaper earns, relative to a partner that is merely slow?
  This is answered by an ablation ladder in which one ingredient changes
  per rung (Section~\ref{sec:ladder}).
\item \textbf{Transfer.} Does a trained shaper, frozen, steer a fresh
  learner differently from a frozen hawk?
  This is the claim of Garcia Segura et al.\ \cite{garciasegura2025}
  stated for this game.
\end{enumerate}
The evidence takes one form throughout: an estimated quantity of
\eqref{eq:estimators}, compared with a value fixed in advance by the
dynamic programme or with the same quantity in a paired control, reported
per seed with its interval and summarised across seeds by sign agreement.
With three to five seeds no test statistic is reported; effect sizes
against exact cells are.

\section{From the deterministic game to stochastic regeneration}
\label{sec:a-to-b}

A shaping effect is measured against a counterfactual partner that does
not shape.
Write \(V_{\mathrm{h}}(\sigma)\) for the return of a partner committed to
take-2 when the learner plays \(\sigma\).
The largest effect a shaper can have on its own return is
\begin{equation}
  \Delta^{*}
  =
  V_{\mathrm{h}}(\sigma_{\mathrm{induced}})-V_{\mathrm{h}}(\sigma_{\mathrm{control}}),
  \label{eq:effect-bound}
\end{equation}
the value of the learner policy it can induce minus the value of the
policy the control partner induces.
\eqref{eq:effect-bound} is a bound on power: no experiment can detect an
effect larger than the game offers.

\paragraph{Stage A.}
After any safe opening the learner's best continuation is a constant
action, so the deterministic game is decided in round~1.
A naive learner leaves opening 2 against any hawk on its own
(Section~\ref{sec:pilot-results}), so \(\sigma_{\mathrm{control}}\) already
delivers the hawk its best outcome: \(V_{\mathrm{h}}=72\) against a
partner that restrains at the opening, and \(72\) against a partner that
has learned the feedback rule (Table~\ref{tab:noise-dp}).
The bound \eqref{eq:effect-bound} is at most one unit.
A null in Stage~A is therefore uninformative about the objective, and a
positive result in Stage~A would be an artefact.
Hypothesis H-A below tests this prediction with the same arms as Stage~B,
which makes the move from A to B a result of the thesis rather than an
assumption of it.

\paragraph{Stage B.}
Under the multiplier \(\xi_t\) of \eqref{eq:growth} the continuation is
a decision (Section~\ref{sec:noise-arm}).
A hawk earns \(72\) against a partner that plays the feedback rule
``1 if \(R<12\), else 2'', \(35.7\) against a partner that restrains only
at the opening, and \(7.4\) against mutual take-2, all exact
(Table~\ref{tab:noise-dp}).
If the naive learner does not acquire the feedback rule on its own, the
bound is \(72-35.7=36.3\) on the hawk's side; if the shaper does more
than commit to take-2 and exploits a partner it has taught to restrain,
the best response to the feedback rule is \(105.5\)
(Table~\ref{tab:dp3}), so the bound is \(69.8\).
Whether the learner acquires the rule alone is check C2 below; the bound
that applies to H-B is conditional on it.
The learner's own stake is of the same size (\(68.8\) against \(34.7\),
survival \(1.00\) against \(0.40\)), so both parties gain from the
feedback rule and the question is whether the shaping objective supplies
what a slow partner does not.

\paragraph{Estimators.}
An epoch is \(n=15\) games.
For agent \(i\), opening leave-2, survival, mean return, and
stock-conditioned action frequency in an epoch are
\begin{equation}
  \begin{aligned}
  \hat{L}_i
  &=
  \frac{1}{n}\sum_{j=1}^{n}\mathbb{1}\bigl[a^{(j)}_{i,1}\ne 2\bigr],
  &\quad
  \hat{S}
  &=
  \frac{1}{n}\sum_{j=1}^{n} S^{(j)},
  \\
  \hat{G}_i
  &=
  \frac{1}{n}\sum_{j=1}^{n} G_i^{(j)},
  &\quad
  \hat{\pi}_i(a\mid R)
  &=
  \frac{\sum_{j,t}\mathbb{1}[R^{(j)}_t=R,\,a^{(j)}_{i,t}=a]}
       {\sum_{j,t}\mathbb{1}[R^{(j)}_t=R]}.
  \end{aligned}
  \label{eq:estimators}
\end{equation}
Last-three-epoch quantities pool \(n=45\).
The tables write those estimators in shorthand a cold reader can decode
without the laboratory:
\begin{itemize}
\item \textbf{Last-open.} Counts of the 15 first-round harvests in
  epoch~15. \(10\times 0+5\times 1\) means ten of those games opened 0
  and five opened 1, not ten epochs of 0.
\item \textbf{Last surv / last ret.} \(\hat{S}\) and \(\hat{G}\) on
  those same 15 games. Stay-on-2 against a hawk returns about 7 and
  dies around round~4; open-1-then-2s returns 71 on the deterministic
  lock.
\item \textbf{Leave-2 last3.} \(\hat{L}\) pooled over epochs 13--15
  (45 openings).
\item \textbf{First majority.} First epoch with \(\hat{L}>8/15\).
\end{itemize}
The opening is round~1 of a 36-round game. Later tokens on a surviving
tape are not last-open.

For a count \(k\) out of \(n\), the \(95\%\) Wilson interval is
\begin{equation}
  \frac{\hat{p}+z^2/2n \;\pm\; z\sqrt{\hat{p}(1-\hat{p})/n+z^2/4n^2}}{1+z^2/n},
  \qquad \hat{p}=k/n,\ z=1.96,
  \label{eq:wilson}
\end{equation}
so \(13/15\) is \([0.62,0.96]\), \(9/15\) is \([0.36,0.80]\) and
\(15/15\) is \([0.80,1.00]\).
Seeds are the unit of replication.
Runs dated 6~September or earlier shared one action-sampling stream
(Chapter~\ref{ch:method}); the 8--9~September packet re-seeds after
agent construction and is three independent seeds on L40S.
Every contrast between configurations is reported as the per-seed
difference and its sign across seeds.
In Stage~B the noise table at a given (epoch, game, round) is identical
across configurations for the same seed, so a between-arm difference on
seed \(s\),
\begin{equation}
  \Delta_s(\hat{Y})=\hat{Y}^{(\mathrm{arm}\,1)}_s-\hat{Y}^{(\mathrm{arm}\,2)}_s,
  \label{eq:paired}
\end{equation}
is a paired comparison in which the environmental draws cancel; with
three seeds, agreement of the sign of \(\Delta_s\) across all seeds is
reported rather than a test statistic.


Per-epoch series of the same quantities are the learning curves; the
window statistics are taken over the last \(W=20\) epochs of a run
(\(W=3\) for the fifteen-epoch pilots).
Two derived quantities describe the dynamics: the first epoch at which a
share exceeds one half (first majority), and stationarity, defined as the
last-window mean differing from the preceding-window mean by less than
the pooled standard error.
Every estimator is implemented in \texttt{cpr\_eval.py} and tested
against records simulated from policies of known value
(\texttt{tests/test\_cpr\_eval.py}) before it is applied to a training
run.

\section{Arms and the ablation ladder}
\label{sec:ladder}

Agent~1 is always a naive learner with the hyperparameters of
Table~\ref{tab:hyperparameters}.
The arms differ in agent~2 (Table~\ref{tab:arms}) and are run in both
stages with the same seeds and, in Stage~B, the same noise tables.

\begin{table}[ht]
\centering
\small
\caption{Arms of the final grid.
  Each rung of the ladder from naive--naive to naive--shaper changes one
  ingredient of agent~2; the configuration files differ in exactly the
  key named in the last column
  (\texttt{scripts/make\_grid\_configs.py}).}
\label{tab:arms}
\begin{tabularx}{\textwidth}{l>{\raggedright\arraybackslash}p{4.6cm}X}
\toprule
Arm & Agent 2 & Isolates (key that changes) \\
\midrule
Test A & frozen always-2 & what a naive learner can learn against a perfectly stationary hawk; the ceiling for C2 \\
naive--naive & naive, learning rate \(1.41\times 10^{-6}\), clip 0.2 & symmetric baseline: who restrains, whether coordination breaks under \(\xi\) \\
slow-LR & naive, learning rate \(3\times 10^{-7}\), clip 0.1 & stationarity of the partner alone (\texttt{learning\_rate}, \texttt{cliprange}) \\
info-off shaper & trial objective \eqref{eq:shaper-objective}, no trial prompt & the trial-level objective (\texttt{is\_shaper}) \\
naive--shaper & trial objective and trial prompt & the trial prompt (\texttt{transmit\_info}) \\
transfer & frozen epoch-100 shaper adapter & whether the trained shaper's policy steers a fresh learner (Stage~B only) \\
\bottomrule
\end{tabularx}
\end{table}

The rung slow-LR to info-off is the test of the objective itself: learning
rate, clip and prompt are equal and only the once-per-trial update with
trial-level advantages differs.
Without it the contrast naive--shaper against slow-LR cannot separate
the objective from the prompt.
The cell ``trial prompt without trial objective'' is not run; if H-B3 is
positive and H-B2 null it becomes the follow-up.

\section{Hypotheses}
\label{sec:hypotheses}

All quantities are last-window estimates of \eqref{eq:estimators};
\(\hat{L}_i(R<12)\) is leave-2 restricted to unmasked rounds with
\(0<R_t<12\); superscripts A and B denote the stage; \(s\) indexes seeds;
\(\Delta_s\) is the paired difference \eqref{eq:paired} between the arms
named.
Two checks precede the hypotheses.
\begin{align}
  \text{C1 (opening invariance):}\quad
  & \hat{L}^{\mathrm{B}}_{1}(s)\in
    \mathrm{CI}_{95}\bigl(\hat{L}^{\mathrm{A}}_{1}(s)\bigr)
    \ \text{ for each } s \text{ (Test~A)},
  \label{eq:c1}\\[3pt]
  \text{C2 (learnability):}\quad
  & \hat{L}^{\mathrm{B}}_{1}(R<12)\ge 0.5
    \ \text{ and }\
    \widehat{\Pr}(S=1\mid a_{1,1}\ne 2)\ge 0.40 .
  \label{eq:c2}
\end{align}
C1 says the noise did not change the opening incentive, which the
dynamic programme predicts (Table~\ref{tab:noise-level}); a shift from
opening 1 to opening 0 within leave-2 is consistent with it, since
\(Q(0)>Q(1)\).
C2, evaluated on Stage~B Test~A, says a naive learner can acquire
stock-conditioned restraint against a stationary hawk within the
horizon; the constants are the
constant-continuation survival cell and a majority of low-stock rounds.
The hypotheses follow.
\begin{align}
  \text{H-A:}\quad
  & \bigl|\Delta_s(\hat{G}_2)\bigr|\le 1
    \ \text{ and }\ \Delta_s(\hat{L}_1)\in\mathrm{CI}_{95}
    \quad\text{for every } s \ \text{(Stage~A)},
  \label{eq:ha}\\[3pt]
  \text{H-B}k\text{:}\quad
  & \Delta^{(k)}_s(\hat{G}_2)>0
    \ \text{ and }\ \Delta^{(k)}_s\bigl(\hat{L}_1(R<12)\bigr)>0
    \quad\text{for every } s,\ k\in\{1,2,3\} \ \text{(Stage~B)},
  \label{eq:hb}\\[3pt]
  \text{H-B4:}\quad
  & \hat{G}^{\mathrm{shaper}}_{2}(s)-1.96\,\widehat{\mathrm{se}}
    > 72 \quad\text{for every } s,
  \label{eq:hb4}\\[3pt]
  \text{H-T:}\quad
  & \Delta^{(\mathrm{T})}_s\bigl(\hat{L}_1(R<12)\bigr)>0
    \ \text{ and }\ \Delta^{(\mathrm{T})}_s(\text{first majority})<0
    \quad\text{for every } s.
  \label{eq:ht}
\end{align}
H-A is the contrast naive--shaper minus slow-LR in Stage~A.
The rungs \(\Delta^{(k)}_s\) are \(k=1\): slow-LR minus naive--naive
(stationarity); \(k=2\): info-off minus slow-LR (trial objective);
\(k=3\): naive--shaper minus info-off (trial prompt); the overall
contrast naive--shaper minus slow-LR is reported alongside.
H-B4 compares the shaper's return with the plain-hawk cell \(72\) and
reads the excess against the ceiling \(105.5\).
\(\Delta^{(\mathrm{T})}_s\) is the fresh learner against the frozen
shaper minus the same learner against the frozen hawk.
Survival and the first-majority epoch of low-stock restraint are reported
for every rung as secondary readouts.

\paragraph{Nulls.}
For H-A the null is the prediction: no rung in Stage~A moves the shaper's
return by more than the one-unit bound.
For H-B\(k\) the null is a mixed sign of \(\Delta_s\) across seeds or a
mean below the sampling interval; the effect size is reported against the
bound of Section~\ref{sec:a-to-b} whichever way it falls.
For H-B4 the null is a shaper return inside the interval of the plain-hawk
cell, which means the shaper committed to take-2 and nothing more.
For H-T the null is no difference between the frozen shaper and the
frozen hawk in what the fresh learner acquires.
A null on any hypothesis is reported as a null with its effect size; none
is rescued by a secondary readout.

\section{Stage B: multiplicative noise on the growth increment}
\label{sec:noise-arm}


The deterministic lock limits what a shaping experiment on it can show.
Against a near-constant partner the fate of the pool is decided by the
joint opening: first-round \((1,2)\) moves stock \(8\to 9\), from which
mutual take-2 is a fixed point; first-round \((2,2)\) leaves the stock at
8 and mutual take-2 collapses around round~4.
Opening 0 versus a hawk moves stock \(8\to 11\), not onto that fixed
point (Chapter~\ref{ch:env}).
After a safe opening, a constant continuation needs no stock information.
A hawk in a hawk--dove split gains one unit (72 against 71) from its
partner's restraint.
Stage~B multiplies the locked growth increment by a mean-one three-point
shock so that the continuation, not only the opening, is a decision
problem.
The deterministic increment is the mean of a per-unit Bernoulli regrowth
process, the scalar form of the density-dependent respawn of Commons
Harvest \cite{perolat2017}; the multiplier is a bounded, mean-preserving
discretisation of that process's spread, chosen so that the constant-pair
cells of Table~\ref{tab:chicken} keep their deterministic values, which
the exact per-unit process does not (Appendix~\ref{app:binomial}).
In the terms of Reed \cite{reed1979} and Sethi et al.\ \cite{sethi2005}
the arm studies growth uncertainty alone, with the shock on the surplus
rather than on the whole recruitment so that collapse remains
attributable to harvest; the joint optimum and the safe feedback rules
below are constant-escapement policies in Reed's sense
(Appendix~\ref{app:reed}).

\subsection{Form}

Let \(R_t\) be the post-harvest stock in round \(t\).
The Chapter~\ref{ch:env} update is
\(R_{t+1}=\min\bigl(K,\, R_t + g_t(R_t)\bigr)\) with
\begin{equation}
  g_t(R)
  \;=\;
  \mathrm{round\_half\_even}\bigl(
    \xi_t\,\rho\, R\,(K-R)/K
  \bigr),
  \qquad
  \xi_t\in\{0.7,\,1.0,\,1.3\}
  \text{ equally likely},
  \label{eq:noise-update}
\end{equation}
\(\rho=0.9\).
The multiplier has mean one, bounded support, and is applied to the
growth increment only.
The post-harvest stock is never perturbed.
In code the multiplier is \(\texttt{xi\_tenths}\in\{7,10,13\}\) inside the
integer Fraction
\(9\,\xi\,R\,(K-R)/(100K)\).
\(\texttt{xi\_tenths}=10\) is bit-identical to Stage~A.
\(\xi\) is never applied at post-harvest \(R=0\).
Draws are independent across rounds and across parallel games, and are
shared by both players within a game.
A discrete support keeps the state space finite (41 stocks, 36 rounds,
three draws), so the game is solved exactly by dynamic programming
(\texttt{cpr\_xi.py}) rather than estimated by simulation.

\subsection{Level and invariance}

The level \(\pm 30\%\) is the largest symmetric three-point multiplier at
which the constant-pair survival cells of Table~\ref{tab:chicken} are
unchanged.
Under hawk versus dove the post-harvest stock is 5 and deterministic
growth is \(3.94\to 4\); with multiplier \(0.7\) the growth is
\(2.76\to 3\), exactly the joint harvest, so the pairing remains safe.
Under mutual take-2 from \(R_0=8\) the post-harvest stock is 4 and
deterministic growth is \(3.24\to 3\); with multiplier \(1.3\) the growth
is \(4.21\to 4\), returning the stock to 8, so the profile never reaches
the fixed point at 9.
(The all-\(\xi=1.3\) sequence is a live 8-cycle; under i.i.d.\ draws that
path has probability \(3^{-36}\).)
At \(\pm 40\%\) both bounds fail: multiplier \(1.4\) lifts mutual take-2
to 9 with positive probability, and multiplier \(0.6\) drives hawk--dove
below 8.
Table~\ref{tab:noise-level} records the i.i.d.\ survival probabilities.

\begin{table}[ht]
\centering
\caption{Survival probability to the horizon for constant action pairs
  and for open-1-then-2s, as a function of the three-point noise level.
  Exact dynamic-programming values from \texttt{cpr\_xi.py}.
  \(\pm 30\%\) leaves the constant-pair cells of Table~\ref{tab:chicken}
  unchanged; larger levels do not.
  The mutual take-2 entry at \(\pm 30\%\) is not exactly zero
  (\(P<10^{-15}\)).}
\label{tab:noise-level}
\begin{tabular}{lcccc}
\toprule
Noise level & \((1,1)\) & \((2,2)\) & \((2,1)\) & open 1 then 2s \\
\midrule
deterministic & 1.00 & 0.00 & 1.00 & 1.00 \\
\(\pm 30\%\)  & 1.00 & 0.00 & 1.00 & 0.80 \\
\(\pm 40\%\)  & 1.00 & 0.15 & 0.83 & 0.78 \\
\(\pm 50\%\)  & 1.00 & 0.15 & 0.84 & 0.67 \\
\bottomrule
\end{tabular}
\end{table}

Invariance of the constant-pair cells is the design property that makes
Stage~B comparable with Stage~A.
The opening decision faces the same constant-pair incentives in both
arms.
Opening values against a frozen always-2 partner, then best-response
continuation, are \(Q=(105,104,102,6)\) in the deterministic game and
\(Q=(102.5, 101.7, 99.1, 91.1)\) under \(\pm 30\%\) noise, for openings
\(0,1,2,3\).
Same ordering on \(0,1,2\); the gap between openings 1 and 2 is \(2.0\)
and \(2.6\) respectively.
Opening 3 then best-response continuation is \(91.1\) under noise because
a stock-conditioned continuer can recover; that is not the untrained
prior.

\subsection{Structure of the stochastic game}

What the noise changes is the continuation.
Under mutual take-2 the stock is a random walk with three regions:
from \(R\ge 11\) the minimum draw returns the stock to at least 11, so
mutual take-2 is safe for the rest of the horizon;
from \(R\le 7\) collapse is certain within a few rounds;
and \(R\in\{8,9,10\}\) is a gamble.
The deterministic fixed point at 9 lies inside the gamble region.
A first \((1,1)\) or \((0,2)\) places the stock in \(\{9,11,12\}\), so
open-0-then-2s against a hawk survives with probability \(0.80\)
(learner return \(58.3\)); a first \((1,2)\) places it in
\(\{8,9,10\}\), so open-1-then-2s against a hawk survives with
probability \(0.40\) (learner return \(34.7\)).
The rule ``take 1 if \(R<12\), else take 2'' survives with probability
one against a frozen hawk (return \(72/68.8\)).
The rule ``take 3 if \(R\ge 14\), else take 0'' survives with probability
one and returns \(102.7\) per player when both use it.
Table~\ref{tab:noise-dp} collects the values.

\begin{table}[ht]
\centering
\caption{Expected return per player and survival probability for
  representative policies in the deterministic game and under
  \(\pm 30\%\) multiplicative noise.
  Exact dynamic-programming values from \texttt{cpr\_xi.py}, rounded to
  one decimal except \((2,2)=7.44\).
  ``Hawk'' is the frozen always-2 partner of Test~A.
  The first three rows are the constant-pair cells and are invariant.
  Every ``good opening, then constant action'' policy becomes risky under
  noise; stock-conditioned feedback rules remain safe.}
\label{tab:noise-dp}
\resizebox{\textwidth}{!}{%
\begin{tabular}{lcccc}
\toprule
 & \multicolumn{2}{c}{Deterministic} & \multicolumn{2}{c}{\(\pm 30\%\) noise} \\
\cmidrule(lr){2-3}\cmidrule(lr){4-5}
Policy pair & Return & Survival & Return & Survival \\
\midrule
\((1,1)\) throughout & 36 / 36 & 1.00 & 36.0 / 36.0 & 1.00 \\
\((2,2)\) throughout & 7 / 7 & 0.00 & 7.44 / 7.44 & 0.00 \\
\((2,1)\) hawk versus dove & 72 / 36 & 1.00 & 72.0 / 36.0 & 1.00 \\
\midrule
\((1,1)\) once, then \((2,2)\) & 71 / 71 & 1.00 & 59.3 / 59.3 & 0.80 \\
Hawk versus ``1 once, then 2'' & 72 / 71 & 1.00 & 35.7 / 34.7 & 0.40 \\
Hawk versus ``0 once, then 2'' & 72 / 70 & 1.00 & 60.3 / 58.3 & 0.80 \\
\((0,0)\) once, then \((3,3)\) & 105 / 105 & 1.00 & 38.7 / 38.7 & 0.25 \\
\midrule
Hawk versus ``1 if \(R<12\), else 2'' & 72 / 69 & 1.00 & 72.0 / 68.8 & 1.00 \\
Both ``2 if \(R\ge 9\), else 1'' & 71 / 71 & 1.00 & 70.8 / 70.8 & 1.00 \\
Both ``3 if \(R\ge 14\), else 0'' & 105 / 105 & 1.00 & 102.7 / 102.7 & 1.00 \\
\midrule
Best response to hawk & 105 & 1.00 & 102.5 & 1.00 \\
Symmetric joint optimum (per player) & 105 & 1.00 & 103.3 & 1.00 \\
\bottomrule
\end{tabular}%
}
\end{table}

Two rows of Table~\ref{tab:noise-dp} define the stakes for the shaping
question.
A hawk that commits to take-2 obtains 72 if its partner uses the
feedback rule and 35.7 if its partner restrains only in round~1; in the
deterministic game the corresponding figures are 72 and 72.
The noise therefore creates a 36-unit interest, on the part of a
committed hawk, in its partner's stock-conditioned restraint, where the
deterministic game offered none, and the partner's cost of providing that
restraint relative to its deterministic return is about two units
(68.8 against 71).
A shaper that did more than commit to take-2 could do better still: the
best response to a partner playing the feedback rule is \(105.5\) under
noise (\(107\) deterministic), so the ceiling on what a shaper gains from
a partner it has taught to restrain is \(105.5\) against the \(72\) of a
plain hawk (Table~\ref{tab:dp3}).
These values are from \texttt{cpr\_xi.py}.


\section{Hyperparameters and sensitivity}
\label{sec:sensitivity}

Table~\ref{tab:hyperparameters} lists every hyperparameter.
Table~\ref{tab:provenance} states where each came from, because several
that earlier drafts described as inherited from ShapeLLM are not: the
released ShapeLLM configurations use an entropy coefficient of zero and
a value coefficient of \(0.1\) for two learners, and a clip of \(0.1\), a
value coefficient of \(0.01\) and a KL coefficient of \(2.0\) with target
\(1\) for a learner against a fixed opponent.

\begin{table}[ht]
\centering
\small
\caption{Provenance of the hyperparameters.
  ``Ours'' means chosen in this project; the sensitivity column names the
  alternative value checked in Section~\ref{sec:sensitivity}.}
\label{tab:provenance}
\begin{tabularx}{\textwidth}{>{\raggedright\arraybackslash}p{2.7cm}>{\raggedright\arraybackslash}p{3.6cm}>{\raggedright\arraybackslash}p{2.6cm}X}
\toprule
Parameter & ShapeLLM release & Live value & Status; sensitivity check \\
\midrule
Naive learning rate & \(1.41\times 10^{-6}\) & \(1.41\times 10^{-6}\) & inherited \\
Shaper learning rate & not released & \(3\times 10^{-7}\) & ours; \(1.41\times 10^{-6}\) \\
Policy clip & 0.2 (two learners), 0.1 (fixed opponent) & 0.2 naive, 0.1 shaper & mixed inheritance \\
Value coefficient \(c_v\) & 0.1 (two learners), 0.01 (fixed opponent) & 0.01 & taken from the fixed-opponent file; 0.1 \\
Entropy coefficient & 0 & \(0.05\to 0.01\) & ours (a pilot at 0.15 bought take-3); 0 \\
KL coefficient & 2.0, target 1 (fixed opponent); default otherwise & 0.2, target 6 (TRL default) & differs; 2.0 with target 1 \\
\(\gamma\), \(\lambda\) & defaults & 1, 0.97 & ours \\
Minibatch & 10 / 6 & 5 & ours \\
Advantage operator & whitening (TRL default) & whitening & inherited; centring in the pilots \\
\bottomrule
\end{tabularx}
\end{table}

A full sweep is neither affordable nor what the argument needs.
The parameters marked ours or mixed are checked one at a time from the
live setting on the cheapest arm that exercises the learner, Stage~B
Test~A, one seed, the full epoch budget: entropy \(0\) against
\(0.05\); \(c_v=0.1\) against \(0.01\); the ShapeLLM fixed-opponent KL
setting against the TRL default; and the shaper learning rate
\(1.41\times 10^{-6}\) against \(3\times 10^{-7}\) on the naive--shaper
arm, because that is the shaper's own parameter and the timescale
confound is this thesis's concern.
The readout is the C2 triple (low-stock restraint, conditional survival,
return).
The result is reported as one table with the statement that the
learnability result is or is not sensitive to each parameter; the grid is
not rerun on that basis unless a setting flips C2, in which case rule R3
below applies.

\section{Protocol}
\label{sec:protocol}

\paragraph{Seeds and platform.}
Three seeds per arm, five on the Stage~B ladder where the thesis claim
lives.
Every run on one NVIDIA L40S; the Apple-silicon pilots are not results,
because bf16 kernels differ between backends and the seed audit of
Chapter~\ref{ch:method} showed platform-dependent divergence.
The entry points re-seed after the trainer is constructed, so seeds are
independent draws; the first epoch of seeds 0 and 1 is compared before
any arm continues and the launch stops if their openings are identical.

\paragraph{Length.}
Every arm runs 100 epochs (1{,}500 games, 500 naive updates, 100 shaper
updates), with the stationarity criterion of Section~\ref{sec:questions}
evaluated over epochs 81--100.
Paired arms share a length: if the naive--shaper arm is not stationary on
any seed, naive--shaper and slow-LR are extended together to 200 on all
seeds.
The Stage~B noise tables are generated at 200 epochs for every seed and
sliced, so an extension keeps the draws of its first 100 epochs and every
arm sees the same draw at the same (epoch, game, round)
(\texttt{tests/test\_cpr\_xi.py}).

\paragraph{Transfer.}
The epoch-100 adapter of each naive--shaper seed is saved; a fresh naive
learner is trained against it, frozen, for 30 epochs, with the shaper's
own prompt rendered as in training, and compared with Stage~B Test~A.

\paragraph{Decision rules.}
R1: Stage~A Test~A with whitened advantages launches first; if any seed's
last-window leave-2 share is below one half, the grid runs with centred
advantages and the operator contrast returns as a result.
R2: the length rule above.
R3: if C2 fails on every seed, H-B is read against the
constant-continuation cell (\(35.7\)) rather than the feedback cell
(\(72\)), and the shaping question becomes whether the shaper makes the
learner acquire what it could not acquire alone.

\paragraph{Budget.}
Measured on the L40S, one two-learner epoch costs about 50 seconds and
one Test~A epoch about 20 seconds.
Stage~B (Test~A, four ladder arms, transfer) at three seeds is about
20 GPU-hours; Stage~A the same without transfer; the sensitivity block
five hours; two further seeds on the Stage~B ladder eight hours.
On five GPUs the whole grid is one working day.

\paragraph{Readouts.}
\texttt{scripts/evaluate\_grid.py} reads the records of every arm and
writes, per stage, the per-arm table of last-window estimates with
intervals, the paired ladder table with sign agreement, \(\hat{\pi}_1(a\mid
R)\) over the last window per arm, and the learning curves; the tables in
Chapter~\ref{ch:results} are those files.

\section{Pilot runs and the decisions they fixed}
\label{sec:pilots}

Fifteen-epoch runs were executed while the design was being fixed:
Tests A and B against frozen bots, the four two-learner arms in Stage~A
and Stage~B, two 50-epoch single-seed extensions, and, after the seed fix,
Test~A under both operators and the Stage~B naive--shaper and slow-LR
arms with independent seeds.
They are reported in Section~\ref{sec:pilot-results} as pilots.
They fixed five decisions.
An entropy coefficient of \(0.15\) moved the learner onto take-3 rather
than restraint; \(0.05\) is used.
Centred advantages left opening 2 within 15 epochs and whitened advantages
did so on a slower, more variable schedule (locking from epoch 23 on a
30-epoch seed); the grid uses the TRL default with rule R1.
The seed audit showed the nominal seeds shared one sampling stream; the
entry points were fixed and the seed-fixed pilots confirmed independent
divergence.
With independent seeds the 15-epoch contrast naive--shaper against
slow-LR had mixed sign on every readout, where the shared-stream pilots
had agreed on every seed; that is why the grid runs 100 epochs and five
seeds on that contrast.
The mixed Apple-silicon and L40S provenance of the early pilots is why the
grid runs on one platform.
